%-*-latex-*-
%-*-latex-*-
\newcommand\COURSE{ciss350}
\newcommand\ASSESSMENT{q10301}
\newcommand\ASSESSMENTTYPE{Quiz}
\newcommand\POINTS{\textwhite{xxx/xxx}}

\input{myquizpreamble}
\input{yliow}
\input{\COURSE}
%-*-latex-*-

%-*-latex-*-

%-*-latex-*-
\renewcommand\TITLE{\ASSESSMENTTYPE \ \ASSESSMENT}

\renewcommand\EMAIL{}
\renewcommand\AUTHOR{eagraham1@cougars.ccis.edu}

\textwidth=6in
\begin{document}
\topmatterthree

All questions below refer to our standard bubblesort.
There are slight variations to the bubblesort algorithm.
You must use the one in my notes/classes.
For big-O,
your answer should be of the form $O(n^k)$ for $k=0, 1, 2, 3, ...$.
Of course $O(n^0)$ is $O(1)$.

%------------------------------------------------------------------------------
\nextq
Write down the array $\{ 5, 3, 1, 4, 2 \}$ after two passes.
\\
\ANSWER
\begin{answerlong}
$\{ 1, 3, 2, 4, 5 \}$
\end{answerlong}

%------------------------------------------------------------------------------
\nextq
For the array $\{ 5, 3, 1, 4, 2 \}$ how many swaps are performed
when the sorting ends?
\\
\ANSWER
\begin{answerlong}
8
\end{answerlong}

%------------------------------------------------------------------------------
\nextq
What array of size 5 that containing the values 1,2,3,4,5 will cause
the algorithm to perform the most number of swaps?
(This will give the highest wallclock time for integer array of size $5$.)
\\
\ANSWER
\begin{answerlong}
$\{ 5, 4, 3, 2, 1 \}$
\end{answerlong}

%------------------------------------------------------------------------------
\nextq
What array of size 5 that contains the values 1,2,3,4,5 will cause
the algorithm to perform the least number of swaps?
(This will give the lowest wallclock time for integer array of size $5$.)
\\
\ANSWER
\begin{answerlong}
$\{ 1, 2, 3, 4, 5 \}$
\end{answerlong}

%------------------------------------------------------------------------------
\nextq
Is it stable? Write YES or NO.
\\
\ANSWER
\begin{answerlong}
YES
\end{answerlong}

%------------------------------------------------------------------------------
\nextq
If it is not stable write down the simplest array
(of smallest size and using smallest positive integers)
to show that it is so.
If it is stable, erase all contents in the answer box and leave it blank.
\\
\ANSWER
\begin{answerlong}
\end{answerlong}

%------------------------------------------------------------------------------
\nextq
What is the big-O of the runtime
in terms of $n$ where $n$ is the size of the array?
\\
\ANSWER
\begin{answerlong}
$O(n^2)$
\end{answerlong}

%------------------------------------------------------------------------------
\nextq
What is the big-O of the best runtime
in terms of $n$ where $n$ is the size of the array?
\\
\ANSWER
\begin{answerlong}
$O(n)$
\end{answerlong}

%------------------------------------------------------------------------------
\nextq
What is the space complexity (i.e., memory usage)
in terms of $n$ where $n$ is the size of the array?
\\
\ANSWER
\begin{answerlong}
$O(1)$
\end{answerlong}

\textsc{Note.}
Note that an array \verb!x! of size $n$ that gives you the highest
the number of swaps (and therefore highest wallclock time)
and an array \verb!y! of size $n$ that gives you the lowers number of
swaps (and therefore lowest wallclock time)
can both give you the same big-O runtime.
Why? Because the first might give you a runtime $T_1(n) = A_1n + B_1$
and the second gives you $T_2(n) = A_2n + B_2$
where $A_1 > A_2$.
But when you compute the big-O, $T_1(n) = O(n)$ and $T_2(n) = O(n)$ --
in terms of big-O, they are the same.
So the concept of \lq\lq best and worst wallclock times" is not the same as the
concept of \lq\lq best and worst (big-O) runtimes".

%------------------------------------------------------------------------------
%-*-latex-*-
\newpage
\input{instructions}
\end{document}

